\documentclass[12pt]{article}

\usepackage[margin=1in]{geometry}
\usepackage{amsmath,amssymb}
\usepackage[T1]{fontenc}
\usepackage{lmodern}

\title{Homework 2}
\author{Hanvesh Aduru}
\date{\today}

\begin{document}

\maketitle

% ============================================================
% PROBLEM 1
% ============================================================

\section*{1. Analyze the Logical Forms}

\subsection*{(a)}

Let
\[
J = \text{John is telling the truth}
\]
and
\[
B = \text{Bill is telling the truth}.
\]

The logical form is
\[
\boxed{(J \land B) \lor (\neg J \land \neg B)}
\]

\subsection*{(b)}

Let
\[
F = \text{I will have fish},
\qquad
C = \text{I will have chicken},
\qquad
M = \text{I will have mashed potatoes}.
\]

The logical form is
\[
\boxed{(F \lor C) \land \neg(F \land M)}
\]

\subsection*{(c)}

The statement ``3 is a common divisor of 6, 9, and 15'' is written as
\[
\boxed{3 \mid 6 \land 3 \mid 9 \land 3 \mid 15}
\]

% ============================================================
% PROBLEM 2
% ============================================================

\section*{2. Well-Formed Formulas}

\subsection*{(a)}

\[
\neg(\neg P \lor \neg\neg R)
\]

\[
\boxed{\text{Well-formed}}
\]

\subsection*{(b)}

\[
\neg(P,Q,\land R)
\]

\[
\boxed{\text{Not well-formed}}
\]

The commas are not logical connectives.

\subsection*{(c)}

\[
P \land \neg P
\]

\[
\boxed{\text{Well-formed}}
\]

\subsection*{(d)}

\[
(P \land Q)(P \lor R)
\]

\[
\boxed{\text{Not well-formed}}
\]

There is no logical connective between the two expressions.

% ============================================================
% PROBLEM 3
% ============================================================

\section*{3. Translate Formulas into English}

Let
\[
T = \text{Taxes will go up}
\]
and
\[
D = \text{The deficit will go up}.
\]

\subsection*{(a)}

\[
T \lor D
\]

Taxes will go up or the deficit will go up.

\subsection*{(b)}

\[
\neg(T \land D) \land \neg(\neg T \land \neg D)
\]

Either taxes will go up or the deficit will go up, but not both.

\subsection*{(c)}

\[
(T \land \neg D) \lor (D \land \neg T)
\]

Either taxes will go up and the deficit will not, or the deficit will go up
and taxes will not.

% ============================================================
% PROBLEM 4
% ============================================================

\section*{4. Premises, Conclusions, and Validity}

\subsection*{(a)}

Let
\[
J = \text{Jane wins the math prize},
\]
\[
P = \text{Pete wins the math prize},
\]
and
\[
C = \text{Pete wins the chemistry prize}.
\]

Premises:
\[
\neg(J \land P)
\]
\[
P \lor C
\]
\[
J
\]

Conclusion:
\[
C
\]

Logical form:
\[
\neg(J \land P),\quad P \lor C,\quad J
\quad \therefore \quad C
\]

The argument is \textbf{valid}.

Since Jane wins the math prize, Pete cannot also win the math prize.
Therefore, Pete must win the chemistry prize.

\subsection*{(b)}

Let
\[
B = \text{The main course is beef},
\]
\[
F = \text{The main course is fish},
\]
\[
P = \text{The vegetable is peas},
\]
and
\[
C = \text{The vegetable is corn}.
\]

Premises:
\[
B \lor F
\]
\[
P \lor C
\]
\[
\neg(F \land C)
\]

Conclusion:
\[
\neg(B \land P)
\]

Logical form:
\[
B \lor F,\quad P \lor C,\quad \neg(F \land C)
\quad \therefore \quad \neg(B \land P)
\]

The argument is \textbf{invalid}.

For example, beef and peas satisfies all three premises while making the
conclusion false.

\subsection*{(c)}

Let
\[
J = \text{John is telling the truth},
\]
\[
B = \text{Bill is telling the truth},
\]
and
\[
S = \text{Sam is telling the truth}.
\]

Premises:
\[
J \lor B
\]
\[
\neg S \lor \neg B
\]

Conclusion:
\[
J \lor \neg S
\]

Logical form:
\[
J \lor B,\quad \neg S \lor \neg B
\quad \therefore \quad J \lor \neg S
\]

The argument is \textbf{valid}.

% ============================================================
% PROBLEM 5
% ============================================================

\section*{5. Truth Tables}

\subsection*{(a)}

\[
\neg[P \land (Q \lor \neg P)]
\]

\begin{center}
\begin{tabular}{c c | c | c | c | c}
$P$ & $Q$ & $\neg P$ & $Q \lor \neg P$
& $P \land (Q \lor \neg P)$
& $\neg[P \land (Q \lor \neg P)]$ \\
\hline
T & T & F & T & T & F \\
T & F & F & F & F & T \\
F & T & T & T & F & T \\
F & F & T & T & F & T
\end{tabular}
\end{center}

\subsection*{(b)}

\[
(P \lor Q) \land (\neg P \lor R)
\]

\begin{center}
\begin{tabular}{c c c | c | c | c | c}
$P$ & $Q$ & $R$ & $P \lor Q$ & $\neg P$
& $\neg P \lor R$ & $(P \lor Q) \land (\neg P \lor R)$ \\
\hline
T & T & T & T & F & T & T \\
T & T & F & T & F & F & F \\
T & F & T & T & F & T & T \\
T & F & F & T & F & F & F \\
F & T & T & T & T & T & T \\
F & T & F & T & T & T & T \\
F & F & T & F & T & T & F \\
F & F & F & F & T & T & F
\end{tabular}
\end{center}

% ============================================================
% PROBLEM 6
% ============================================================

\section*{6. Exclusive OR}

The symbol $+$ means exclusive OR:
\[
P + Q = \text{P or Q, but not both}.
\]

\subsection*{(a)}

\begin{center}
\begin{tabular}{c c | c}
$P$ & $Q$ & $P+Q$ \\
\hline
T & T & F \\
T & F & T \\
F & T & T \\
F & F & F
\end{tabular}
\end{center}

\subsection*{(b)}

An equivalent formula using only $\land$, $\lor$, and $\neg$ is

\[
\boxed{(P \lor Q) \land \neg(P \land Q)}
\]

Truth table:

\begin{center}
\begin{tabular}{c c | c | c | c | c}
$P$ & $Q$ & $P \lor Q$ & $P \land Q$
& $\neg(P \land Q)$
& $(P \lor Q) \land \neg(P \land Q)$ \\
\hline
T & T & T & T & F & F \\
T & F & T & F & T & T \\
F & T & T & F & T & T \\
F & F & F & F & T & F
\end{tabular}
\end{center}

The final column is identical to the truth table for $P+Q$.

Therefore,
\[
\boxed{P+Q \equiv (P \lor Q) \land \neg(P \land Q)}
\]

% ============================================================
% PROBLEM 7
% ============================================================

\section*{7. Formula Using Only $\land$ and $\neg$}

Find a formula equivalent to $P \lor Q$ using only $\land$ and $\neg$.

By De Morgan's Law,
\[
\boxed{P \lor Q \equiv \neg(\neg P \land \neg Q)}
\]

Truth table:

\begin{center}
\begin{tabular}{c c | c | c | c | c}
$P$ & $Q$ & $\neg P$ & $\neg Q$
& $\neg P \land \neg Q$
& $\neg(\neg P \land \neg Q)$ \\
\hline
T & T & F & F & F & T \\
T & F & F & T & F & T \\
F & T & T & F & F & T \\
F & F & T & T & T & F
\end{tabular}
\end{center}

The final column is identical to $P \lor Q$.

% ============================================================
% PROBLEM 8
% ============================================================

\section*{8. NAND}

The symbol $|$ means NAND:
\[
P|Q = \neg(P \land Q)
\]

\subsection*{(a)}

\begin{center}
\begin{tabular}{c c | c | c}
$P$ & $Q$ & $P \land Q$ & $P|Q$ \\
\hline
T & T & T & F \\
T & F & F & T \\
F & T & F & T \\
F & F & F & T
\end{tabular}
\end{center}

\subsection*{(b)}

\[
\boxed{P|Q \equiv \neg(P \land Q)}
\]

\subsection*{(c)}

Using only the connective $|$:

For $\neg P$:
\[
\boxed{\neg P \equiv P|P}
\]

For $P \lor Q$:
\[
\boxed{P \lor Q \equiv (P|P)|(Q|Q)}
\]

For $P \land Q$:
\[
\boxed{P \land Q \equiv (P|Q)|(P|Q)}
\]

% ============================================================
% PROBLEM 9
% ============================================================

\section*{9. Formula Equivalence}

The formulas are:

\[
\begin{aligned}
(a) &\quad (P \land Q) \lor (\neg P \land \neg Q)\\
(b) &\quad \neg P \lor Q\\
(c) &\quad (P \lor \neg Q) \land (Q \lor \neg P)\\
(d) &\quad \neg(P \lor Q)\\
(e) &\quad (Q \land P) \lor \neg P
\end{aligned}
\]

Truth table:

\begin{center}
\begin{tabular}{c c | c c c c c}
$P$ & $Q$ & (a) & (b) & (c) & (d) & (e) \\
\hline
T & T & T & T & T & F & T \\
T & F & F & F & F & F & F \\
F & T & F & T & F & F & T \\
F & F & T & T & T & T & T
\end{tabular}
\end{center}

Comparing the final columns:

\[
\boxed{(a) \equiv (c)}
\]

and

\[
\boxed{(b) \equiv (e)}
\]

Formula (d) is not equivalent to any of the other formulas.

\end{document}

